\documentclass[12pt]{article}
\usepackage{amsmath,amsthm,amssymb,mathtools}
\usepackage{hyperref}

\newtheorem{theorem}{Theorem}[section]
\newtheorem{lemma}[theorem]{Lemma}
\newtheorem{proposition}[theorem]{Proposition}
\newtheorem{corollary}[theorem]{Corollary}
\newtheorem{definition}[theorem]{Definition}
\newtheorem{remark}[theorem]{Remark}

\DeclareMathOperator{\diag}{diag}
\DeclareMathOperator{\Res}{Res}
\DeclareMathOperator{\Ad}{Ad}
\newcommand{\A}{\mathbb{A}}
\newcommand{\Z}{\mathbb{Z}}
\newcommand{\Q}{\mathbb{Q}}
\newcommand{\C}{\mathbb{C}}
\newcommand{\calO}{\mathcal{O}}
\newcommand{\calW}{\mathcal{W}}

\title{Foundations: Local Euler Factor Factorization and the Clebsch--Gordan Cancellation\\
\large Proof Document 01 --- sym2-effective-lbound}
\author{}
\date{}

\begin{document}
\maketitle

\begin{abstract}
We prove Theorem~F-1: for any unramified finite place $v$ of a number field $F$,
with Satake parameters $(\alpha_v, \beta_v)$ satisfying $\alpha_v \beta_v = 1$,
the unramified local Jacquet--Shalika integral equals
\[
  I_v(s) = \frac{1}{(1-q_v^{-s})^2(1-\alpha_v^2 q_v^{-s})(1-\beta_v^2 q_v^{-s})}
          = \zeta_v(s) \cdot L_v(s, \pi_v, \Ad).
\]
The key step is an exact cancellation of the $(1+q_v^{-s})$ factor, which follows
from the $\mathrm{SU}(2)$ Clebsch--Gordan identity $\chi_l^2 = \sum_{j=0}^l \chi_{2j}$.
This result is [THM F-1] in the project specification.
\end{abstract}

\tableofcontents

\section{Setup and Notation}

Let $F$ be a number field and $v$ a non-Archimedean place with residue field
$\mathbf{F}_{q_v}$, ring of integers $\calO_v$, and uniformizer $\varpi_v$.
Let $\pi_v$ be an unramified irreducible admissible representation of $GL_2(F_v)$
with trivial central character and Satake parameters $(\alpha_v, \beta_v)$ satisfying
\begin{equation}\label{eq:satake-unit}
  \alpha_v \beta_v = 1.
\end{equation}

The local unramified Whittaker function $W_{f,v}^\circ \in \calW(\pi_v, \psi_v)$
is normalized so that $W_{f,v}^\circ(I_2) = 1$.

\section{Casselman--Shalika Formula}

\begin{theorem}[Casselman--Shalika {\cite[Theorem 5.4, p.~227]{CS80};
Jacquet--Shalika {\cite[Section 2.2, p.~511]{JS81}}}]\label{thm:CS}
For $l \geq 0$,
\[
  W_{f,v}^\circ\bigl(\diag(\varpi_v^l, 1)\bigr)
  = q_v^{-l/2}\, \chi_l(\alpha_v, \beta_v),
\]
where
\[
  \chi_l(\alpha, \beta)
  = \frac{\alpha^{l+1} - \beta^{l+1}}{\alpha - \beta}
\]
is the character of the $(l+1)$-dimensional irreducible representation of $\mathrm{SU}(2)$.
For $l < 0$, $W_{f,v}^\circ(\diag(\varpi_v^l, 1)) = 0$.
\end{theorem}

\begin{proof}[Source-to-normalization bridge]
Casselman--Shalika's Theorem 5.4 evaluates the normalized spherical Whittaker
function on the dominant torus by the Weyl-character sum.  Jacquet--Shalika
transcribe the equivalent essential-vector
normalization
\[
  W(a_J)=\delta^{1/2}(a_J)\operatorname{Tr}\rho_J(A)
\]
for $GL_r$.  Taking $r=2$, $J=(l,0)$, and $a_J=\diag(\varpi_v^l,1)$ gives
$\delta^{1/2}(a_J)=q_v^{-l/2}$; the restriction to $SL_2$ of the
highest-weight representation $\rho_J$ is its $l$-th symmetric power, whose
character is $\chi_l$.  The asserted formula follows under the repository's
measure normalization.
\end{proof}

\section{Iwasawa Reduction of the Local Integral}

The local Jacquet--Shalika integral with test function $\Phi_v = \mathbf{1}_{\calO_v^2}$ is
\[
  I_v(s) = \int_{N_2(F_v) \backslash GL_2(F_v)}
            W_{f,v}(g)\, W_{\tilde f,v}(g)\, \Phi_v(e_2 g)\, |\det g|_v^s\, d^\times g,
\]
where $e_2 = (0,1)$ is the second standard basis row vector.

\begin{lemma}[Iwasawa reduction]\label{lem:iwasawa}
Using the Iwasawa decomposition $GL_2(F_v) = N_2(F_v) A_2(F_v) K_v$ and the change of
variables $a = \diag(\varpi_v^{k}, \varpi_v^{m})$ with $l = k - m$, the support
conditions force $l \geq 0$ and $m \geq 0$, and
\[
  I_v(s) = \underbrace{\left(\sum_{m=0}^\infty q_v^{-2ms}\right)}_{=: S_m}
            \cdot
            \underbrace{\left(\sum_{l=0}^\infty
              \chi_l(\alpha_v, \beta_v)^2\, q_v^{-ls}\right)}_{=: S_l}.
\]
\end{lemma}

\begin{proof}
Standard Iwasawa integration: the modular character and Haar measure yield the
normalization $d^\times a = q_v^l\, dk\, dm$ in the $(l, m)$ variables.  The test
function $\Phi_v(e_2 a) = \mathbf{1}_{\calO_v}(\varpi_v^m)$ restricts to $m \geq 0$.
The Whittaker pairing $W_{f,v}(a) W_{\tilde f,v}(a)$ evaluates via
Theorem~\ref{thm:CS} and the unitarity $W_{\tilde f,v} = \overline{W_{f,v}}$, giving
$\chi_l(\alpha_v,\beta_v)^2$ after noting $\alpha_v \beta_v = 1$ implies
$\overline{\chi_l} = \chi_l$ on the unit circle.
\end{proof}

\section{Evaluation of the Two Series}

\subsection{The $m$-series}

\begin{lemma}\label{lem:m-series}
$S_m = \dfrac{1}{(1-q_v^{-s})(1+q_v^{-s})}$.
\end{lemma}

\begin{proof}
Geometric series: $S_m = \sum_{m=0}^\infty q_v^{-2ms} = (1 - q_v^{-2s})^{-1}
= [(1-q_v^{-s})(1+q_v^{-s})]^{-1}$.
\end{proof}

\subsection{The $l$-series via the Clebsch--Gordan Identity}

The fundamental identity we use is the $\mathrm{SU}(2)$ Clebsch--Gordan decomposition:
\begin{equation}\label{eq:CG}
  \chi_l^2 = \chi_l \cdot \chi_l = \sum_{j=0}^l \chi_{2j}.
\end{equation}

\begin{lemma}[Even-character generating function]\label{lem:even-gen}
With $x = q_v^{-s}$ and $\alpha_v \beta_v = 1$,
\[
  \sum_{j=0}^\infty \chi_{2j}(\alpha_v, \beta_v)\, x^j
  = \frac{1+x}{(1-\alpha_v^2 x)(1-\beta_v^2 x)}.
\]
\end{lemma}

\begin{proof}
Expand:
\[
  \sum_{j=0}^\infty \chi_{2j}(\alpha,\beta) x^j
  = \frac{1}{\alpha-\beta}\left(
      \frac{\alpha}{1-\alpha^2 x} - \frac{\beta}{1-\beta^2 x}\right).
\]
Computing the numerator of the combined fraction:
\[
  \alpha(1-\beta^2 x) - \beta(1-\alpha^2 x)
  = (\alpha-\beta) + (\alpha^2 \beta - \alpha \beta^2) x
  = (\alpha-\beta)(1 + \alpha\beta x)
  = (\alpha-\beta)(1+x)
\]
using $\alpha\beta = 1$.  Dividing by $(\alpha-\beta)(1-\alpha^2 x)(1-\beta^2 x)$
gives the claim.
\end{proof}

\begin{lemma}[$l$-series]\label{lem:l-series}
With $x = q_v^{-s}$,
\[
  S_l = \sum_{l=0}^\infty \chi_l(\alpha_v,\beta_v)^2\, x^l
  = \frac{1+x}{(1-x)(1-\alpha_v^2 x)(1-\beta_v^2 x)}.
\]
\end{lemma}

\begin{proof}
Apply \eqref{eq:CG} and swap the order of summation:
\[
  S_l = \sum_{l=0}^\infty \sum_{j=0}^l \chi_{2j} x^l
  = \sum_{j=0}^\infty \chi_{2j} \sum_{l=j}^\infty x^l
  = \frac{1}{1-x} \sum_{j=0}^\infty \chi_{2j} x^j.
\]
Apply Lemma~\ref{lem:even-gen} to obtain the result.
\end{proof}

\section{Exact Cancellation and Main Theorem}

\begin{theorem}[F-1: Local Euler Factor Factorization]\label{thm:F1}
Under the hypotheses of Section~1,
\[
  I_v(s)
  = \frac{1}{(1-q_v^{-s})^2(1-\alpha_v^2 q_v^{-s})(1-\beta_v^2 q_v^{-s})}
  = \zeta_v(s) \cdot L_v(s, \pi_v, \Ad).
\]
\end{theorem}

\begin{proof}
By Lemma~\ref{lem:iwasawa}, $I_v(s) = S_m \cdot S_l$.  Multiplying
Lemmas~\ref{lem:m-series} and \ref{lem:l-series}:
\[
  I_v(s)
  = \frac{1}{(1-q_v^{-s})(1+q_v^{-s})}
    \cdot
    \frac{1+q_v^{-s}}{(1-q_v^{-s})(1-\alpha_v^2 q_v^{-s})(1-\beta_v^2 q_v^{-s})}.
\]
The factor $(1+q_v^{-s})$ cancels exactly, yielding
\[
  I_v(s)
  = \frac{1}{(1-q_v^{-s})^2(1-\alpha_v^2 q_v^{-s})(1-\beta_v^2 q_v^{-s})}.
\]
The identification with $\zeta_v(s) \cdot L_v(s,\pi_v,\Ad)$ follows from the
dual group decomposition $\mathbf{1} \oplus \Ad$ of $\mathrm{std} \otimes \mathrm{std}^\vee$
for $GL_2$: the adjoint representation has Satake parameters $(\alpha_v^2, 1, \beta_v^2)$
(for $\alpha_v \beta_v = 1$), so
\[
  L_v(s,\pi_v,\Ad)
  = (1-q_v^{-s})^{-1}(1-\alpha_v^2 q_v^{-s})^{-1}(1-\beta_v^2 q_v^{-s})^{-1},
\]
and $\zeta_v(s) = (1-q_v^{-s})^{-1}$.
\end{proof}

\begin{remark}
The cancellation is purely algebraic: it does not depend on any analytic properties
of $\pi_v$ beyond the Casselman--Shalika formula and $\alpha_v \beta_v = 1$.
When $\omega_{\pi,v}$ is non-trivial, the Satake parameters of $\tilde\pi_v$ become
$(\alpha_v^{-1}\omega_v(\varpi_v)^{-1}, \beta_v^{-1}\omega_v(\varpi_v)^{-1})$,
and the result generalizes to $L_v(s,\pi_v,\Ad \otimes \omega_\pi)$.
\end{remark}

\section{Numerical Verification}

The file \texttt{src/euler\_factors.py} implements the computation of $I_v(s)$ via:
\begin{enumerate}
  \item Direct summation of the double series (truncated, for validation)
  \item The closed-form formula from Theorem~\ref{thm:F1}
\end{enumerate}
For specific Satake parameters (e.g., those of the Ramanujan $\Delta$ form at small
primes), both computations agree to machine precision.  The unit tests in
\texttt{tests/test\_euler\_factors.py} verify the $(1+q_v^{-s})$ cancellation
numerically and confirm the closed-form identity.

\begin{thebibliography}{9}
\bibitem{CS80} W.~Casselman and J.~Shalika, The unramified principal series of
  $p$-adic groups II: the Whittaker function, Compositio Math.\ \textbf{41} (1980), 207--231.
\bibitem{JS81} H.~Jacquet and J.~A.~Shalika, On Euler products and the classification of
  automorphic forms I, Amer.\ J.\ Math.\ \textbf{103} (1981), 499--558.
\end{thebibliography}

\end{document}
